% GENERATED FILE. Edit canonical lessons, then run npm run generate:books.
% canonical-source-hash: fnv1a64:29264f822067ab29
% canonical-lessons: LA-C474-aro, LA-C474-sulcus, LA-C474-vomer, LA-C474-arator, LA-C474-occo

\chapter{To plough, Furrow, Ploughshare, Ploughman, To harrow}
\label{ch:la-to-plough-furrow-ploughshare-ploughman-to-harrow}
\hlchaptermodality{474}

\begin{chapteropening}
\textbf{By the end of this chapter:} \emph{I can talk about ploughing: ploughing a field, the furrow, the ploughshare, the ploughman and harrowing.}

Complete the last lesson of chapter 474: I can talk about ploughing: ploughing a field, the furrow, the ploughshare, the ploughman and harrowing.
\end{chapteropening}

\section[arō, arāre]{arō, arāre — to plough}

\label{lesson:LA-C474-aro}



\begin{quote}
Before the new one: say the Latin for a cottage, then the Latin for a hut.
\end{quote}

\subsection*{You'll want to know: arō, arāre}
\textbf{arō, arāre} — \textquotedblleft{}to plough\textquotedblright{}.

\begin{quote}
Vēre agricolae agrōs arant.
\end{quote}

\begin{quote}
\emph{In spring the farmers plough the fields.}
\end{quote}

\emph{arō, arāre, arāvī, arātum}. The plough is \emph{arātrum}. English \emph{arable} comes from it.

\begin{cousinweb}
\emph{arō} is inherited from Proto-Italic \emph{araō} (EtymDB), with Greek \emph{aroō} 'to plough' as a cousin (EtymDB; Lewis \& Short); Lewis \& Short add Gothic \emph{arjan}. English \emph{arable} comes from it.
\end{cousinweb}

\begin{grammarlens}[title={What you've built}]
One more word for telling, explaining and describing at B1.
\end{grammarlens}

\subsection*{Your turn}
\begin{itemize}
\raggedright
  \item \emph{Say it:} \emph{arō, arāre}
  \item \emph{Say it:} \emph{arō, arāre}, once more
  \item \emph{Say it:} a sentence of your own with \emph{arō, arāre}
  \item \emph{Recall:} say the Latin for a cottage, then the Latin for a hut, then say \emph{arō, arāre} again
\end{itemize}

\begin{tcolorbox}[breakable,colback=teal!4,colframe=teal!35!black,title={Before you move on}]
What does arō, arāre mean? (\textquotedblleft{}To plough\textquotedblright{}.) Say it once more.
\end{tcolorbox}
\section[sulcus, sulcī]{sulcus, sulcī — a furrow; a trench}

\label{lesson:LA-C474-sulcus}



\begin{quote}
Before the new one: say the Latin for a hut, then the Latin for to plough.
\end{quote}

\subsection*{You'll want to know: sulcus, sulcī}
\textbf{sulcus, sulcī} — \textquotedblleft{}a furrow; a trench\textquotedblright{}.

\begin{quote}
Agricola sulcōs rēctōs in agrō facit.
\end{quote}

\begin{quote}
\emph{The farmer makes straight furrows in the field.}
\end{quote}

Second declension, masculine: \emph{sulcus, sulcī m.} Also the wake of a ship. English keeps \emph{sulcus} in anatomy for a groove.

\begin{cousinweb}
\emph{sulcus} is inherited from PIE \emph{solk-o-} 'furrow' (EtymDB); Lewis \& Short compare Greek \emph{holkos} 'furrow'.
\end{cousinweb}

\begin{grammarlens}[title={What you've built}]
One more word for telling, explaining and describing at B1.
\end{grammarlens}

\subsection*{Your turn}
\begin{itemize}
\raggedright
  \item \emph{Say it:} \emph{sulcus, sulcī}
  \item \emph{Say it:} \emph{sulcus, sulcī}, once more
  \item \emph{Say it:} a sentence of your own with \emph{sulcus, sulcī}
  \item \emph{Recall:} say the Latin for a hut, then the Latin for to plough, then say \emph{sulcus, sulcī} again
\end{itemize}

\begin{tcolorbox}[breakable,colback=teal!4,colframe=teal!35!black,title={Before you move on}]
What does sulcus, sulcī mean? (\textquotedblleft{}A furrow; a trench\textquotedblright{}.) Say it once more.
\end{tcolorbox}
\section[vōmer, vōmeris]{vōmer, vōmeris — a ploughshare}

\label{lesson:LA-C474-vomer}



\begin{quote}
Before the new one: say the Latin for to plough, then the Latin for a furrow.
\end{quote}

\subsection*{You'll want to know: vōmer, vōmeris}
\textbf{vōmer, vōmeris} — \textquotedblleft{}a ploughshare\textquotedblright{}.

\begin{quote}
Vōmer ex ferrō est et terram dūram frangit.
\end{quote}

\begin{quote}
\emph{The ploughshare is made of iron and breaks the hard ground.}
\end{quote}

Third declension, masculine: \emph{vōmer, vōmeris m.} (also \emph{vōmis}). The iron blade of the plough.

\begin{cousinweb}
The origin of \emph{vōmer} is not settled: Lewis \& Short give the sense but no source, and EtymDB lists it without a parent. English keeps \emph{vomer} for a bone of the nose.
\end{cousinweb}

\begin{grammarlens}[title={What you've built}]
One more word for telling, explaining and describing at B1.
\end{grammarlens}

\subsection*{Your turn}
\begin{itemize}
\raggedright
  \item \emph{Say it:} \emph{vōmer, vōmeris}
  \item \emph{Say it:} \emph{vōmer, vōmeris}, once more
  \item \emph{Say it:} a sentence of your own with \emph{vōmer, vōmeris}
  \item \emph{Recall:} say the Latin for to plough, then the Latin for a furrow, then say \emph{vōmer, vōmeris} again
\end{itemize}

\begin{tcolorbox}[breakable,colback=teal!4,colframe=teal!35!black,title={Before you move on}]
What does vōmer, vōmeris mean? (\textquotedblleft{}A ploughshare\textquotedblright{}.) Say it once more.
\end{tcolorbox}
\section[arātor, arātōris]{arātor, arātōris — a ploughman, a farmer}

\label{lesson:LA-C474-arator}



\begin{quote}
Before the new one: say the Latin for a furrow, then the Latin for a ploughshare.
\end{quote}

\subsection*{You'll want to know: arātor, arātōris}
\textbf{arātor, arātōris} — \textquotedblleft{}a ploughman, a farmer\textquotedblright{}.

\begin{quote}
Arātor bovēs ad agrum dūcit.
\end{quote}

\begin{quote}
\emph{The ploughman leads the oxen to the field.}
\end{quote}

Third declension, masculine: \emph{arātor, arātōris m.}, from \emph{arō} with \emph{-tor}.

\begin{cousinweb}
\emph{arātor} is built on \emph{arō} 'to plough' with the suffix \emph{-tor} 'one who' (EtymDB; Lewis \& Short).
\end{cousinweb}

\begin{grammarlens}[title={What you've built}]
One more word for telling, explaining and describing at B1.
\end{grammarlens}

\subsection*{Your turn}
\begin{itemize}
\raggedright
  \item \emph{Say it:} \emph{arātor, arātōris}
  \item \emph{Say it:} \emph{arātor, arātōris}, once more
  \item \emph{Say it:} a sentence of your own with \emph{arātor, arātōris}
  \item \emph{Recall:} say the Latin for a furrow, then the Latin for a ploughshare, then say \emph{arātor, arātōris} again
\end{itemize}

\begin{tcolorbox}[breakable,colback=teal!4,colframe=teal!35!black,title={Before you move on}]
What does arātor, arātōris mean? (\textquotedblleft{}A ploughman, a farmer\textquotedblright{}.) Say it once more.
\end{tcolorbox}
\section[occō, occāre]{occō, occāre — to harrow, to break up the soil}

\label{lesson:LA-C474-occo}



\begin{quote}
Before the new one: say the Latin for a ploughshare, then the Latin for a ploughman.
\end{quote}

\subsection*{You'll want to know: occō, occāre}
\textbf{occō, occāre} — \textquotedblleft{}to harrow, to break up the soil\textquotedblright{}.

\begin{quote}
Agricola agrum prīmum arat, deinde occat.
\end{quote}

\begin{quote}
\emph{The farmer first ploughs the field, then harrows it.}
\end{quote}

\emph{occō, occāre, occāvī, occātum}, from \emph{occa} 'a harrow'.

\begin{cousinweb}
\emph{occō} is built on \emph{occa} 'a harrow' (EtymDB), from the PIE root \emph{h\textsubscript{2}e\'{k}-} 'sharp'. Lewis \& Short give German \emph{eggen} 'to harrow' as a cousin.
\end{cousinweb}

\begin{grammarlens}[title={What you've built}]
One more word for telling, explaining and describing at B1.
\end{grammarlens}

\subsection*{Your turn}
\begin{itemize}
\raggedright
  \item \emph{Say it:} \emph{occō, occāre}
  \item \emph{Say it:} \emph{occō, occāre}, once more
  \item \emph{Say it:} a sentence of your own with \emph{occō, occāre}
  \item \emph{Recall:} say the Latin for a ploughshare, then the Latin for a ploughman, then say \emph{occō, occāre} again
\end{itemize}

\begin{tcolorbox}[breakable,colback=teal!4,colframe=teal!35!black,title={Before you move on}]
What does occō, occāre mean? (\textquotedblleft{}To harrow, to break up the soil\textquotedblright{}.) Say it once more.
\end{tcolorbox}
